\documentclass[10pt,a4paper]{amsart}
\usepackage[margin=19mm]{geometry}
\usepackage{amsmath,amssymb,mathtools}
\usepackage[scr=rsfs,cal=euler]{mathalfa}
\usepackage{xfrac}
\usepackage[unicode,hidelinks,bookmarks=false]{hyperref}
\newcommand{\ie}{i.e.\ }
\newcommand{\cf}{cf.\ }
\newcommand{\resp}{respectively\ }
\newcommand{\Subsection}{\subsection}
\providecommand{\nobreakdash}{\nobreak\hbox{-}}
\setlength{\emergencystretch}{2em}
\multlinegap0pt
\usepackage{tikz}
\usepackage{mathtools}
\usepackage{amssymb}
\newtheorem{thm}{Theorem}
\newtheorem{prop}{Proposition}
\newtheorem{pf}{Pf}
\usepackage{cleveref}
\crefname{thm}{Theorem}{Theorems}
\crefname{prop}{Proposition}{Propositions}
\crefname{pf}{Pf}{Pfs}
\DeclareMathOperator{\PGam}{\mathrm{P}\Gamma}
\DeclareMathOperator{\PU}{PU}
\DeclareMathOperator{\SR}{\mathcal{SR}}
\newcommand{\bbB}{\mathbb{B}}
\newcommand{\bs}{\backslash}
\newcommand{\calC}{\mathcal{C}}
\newcommand{\calD}{\mathcal{D}}
\newcommand{\calH}{\mathcal{H}}
\newcommand{\gam}{\gamma}
\newcommand{\ssm}{\smallsetminus}
\title{Automorphisms of the moduli space of smooth cubic surfaces and its fundamental group}
\author{Gregorio Baldi, Benson Farb, Ariyan Javanpeykar, Matthew Stover}
\date{Source: arXiv:2605.16658v1 (CC BY 4.0)}
\begin{document}
\maketitle

\noindent\emph{Source and licence.} Automorphisms of the moduli space of smooth cubic surfaces and its fundamental group. Version arXiv:2605.16658v1 (CC BY 4.0), distributed by arXiv under the Creative Commons Attribution 4.0 International licence, http://creativecommons.org/licenses/by/4.0/, as recorded in the arXiv licence metadata for this exact version and shown on the abstract page https://arxiv.org/abs/2605.16658v1. Full licence notice: \texttt{licenses/arxiv-2605.16658-CC-BY-4.0.txt}.\par\smallskip

\noindent\textit{Editorial scope.} Counted unit: the article's prop prop:ConnectedSupport (source0.tex lines 625--627). The article's complete original proof of it is delivered with it (source0.tex lines 630--646, one \texttt{proof} environment of 17 lines). No mathematics is written, repaired or re-derived: the statement and the proof are the article's own lines, in the article's own notation, and no retained line is substituted or re-flowed.  Retained uncounted as the source-local context the proof needs: lines 424--424; lines 560--571; lines 583--585; lines 588--610.  Nothing was omitted from any retained span, so the package carries no omission record.  The retained lines cite 1 works, whose entries are transcribed verbatim from the article's own bibliography source in the e-print and delivered with short editorial labels recorded in the item's reference mapping.  The retained lines contain the article's own diagram code, which the wrapper typesets with the standard tikz packages; no image file is delivered and the package declares no asset dependency.  Editorial formatting only, in addition: an AMS \texttt{amsart} preamble that restates the article's own declarations and macros used by the retained lines, namely the environments \texttt{thm}, \texttt{prop}, \texttt{pf} on separate counters, together with the standard packages those lines need.  The retained environments print in this document as Theorem 1, Proposition 1, Pf 1 in order of appearance, whereas the article prints the same items with the numbering of its own full document.  The complete native source of the article as distributed in the arXiv e-print for this exact version is bundled unchanged as \texttt{sources/200735/source0.tex}.
\section{$\calC$ as a divisor complement in a Picard modular $4$-fold}\label{sec:Prelims}

\begin{thm}[Allcock--Carlson--Toledo \cite{ACT}]\label{thm:ACTbig}
With the notation established in this section:
\begin{enumerate}

	\item The image of the period map $\calC \to \PGam_4 \bs \bbB^4$ is precisely $X \ssm D$.

	\item If $r_1, r_2 \in \SR$ then $\bbB^3_{r_1}$ and $\bbB^3_{r_2}$ are either disjoint or meet orthogonally.

	\item The action of $\PGam_4$ on the set of lines spanned by short roots is transitive.

\end{enumerate}
\end{thm}

\begin{thm}\label{thm:ACTgensE6}
Let $R_j \in \PGam_4$ be the images in $\PU(4,1)$ of hexaflections through the roots $r_j$ given in \cite[\S 7]{ACT}. Then the $R_j$ satisfy the relations of $\mathrm{A}(\widetilde{\mathrm{E}}_6)$ on the affine $\mathrm{E}_6$ diagram, where the circled vertices in \Cref{fig:ArtinGens} indicate hexaflections through mutually intersecting hyperplanes associated with short roots. In particular, there is a surjective homomorphism $\mathrm{A}(\widetilde{\mathrm{E}}_6)\to \PGam_4$ factoring through the representation $\pi_1(\calC) \to \PGam_4$.
\end{thm}

%%%%%%%%%%%%%%%%%%%%
\begin{figure}[h]
\centering
\scalebox{0.75}{
\begin{tikzpicture}
\draw[fill=black] (0,0) circle (0.075cm) node [yshift=0.5cm] {$r_3$};
\draw (0,0) circle (0.2cm);
\draw[fill=black] (2,0) circle (0.075cm) node [yshift=0.5cm] {$r_4$};
\draw[fill=black] (-2,0) circle (0.075cm) node [yshift=0.5cm] {$r_2$};
\draw[fill=black] (4,0) circle (0.075cm) node [yshift=0.5cm] {$r_5$};
\draw (4,0) circle (0.2cm);
\draw[fill=black] (-4,0) circle (0.075cm) node [yshift=0.5cm] {$r_1$};
\draw (-4,0) circle (0.2cm);
\draw[fill=black] (0,-2) circle (0.075cm) node [xshift=0.5cm] {$r_6$};
\draw[fill=black] (0,-4) circle (0.075cm) node [xshift=0.5cm] {$r_7$};
\draw (0,-4) circle (0.2cm);
\draw (-4,0) -- (4,0);
\draw (0,0) -- (0,-4);
\end{tikzpicture}
}
\caption{The hexaflection generators for $\pi_1(\calC)$ from $\mathrm{A}(\widetilde{\mathrm{E}}_6)$.}\label{fig:ArtinGens}
\end{figure}
%%%%%%%%%%%%%%%%%%%%

\begin{prop}\label{prop:ConnectedSupport}
The arrangement $\calH$ in $\bbB^4$ defined in \Cref{sec:Prelims} has connected support.
\end{prop}

\begin{pf}
Consider the hyperplanes $\calD_k = \bbB^4_{r_{2 k - 1}}$ with $r_{2 k - 1}$ as in \Cref{thm:ACTgensE6}, $k = 1,2,3$, and set
	\[
	\calH_0 \coloneqq \bigcup_{k = 1,2,3} \calD_k.
	\]
\Cref{thm:ACTbig} implies that the hyperplanes in the support of $\calH_0$ meet orthogonally, since the hexaflections $R_j$ all commute. Moreover, each generator $R_j$ of $\PGam_4$, $1 \le j \le 7$, has the property that $R_j(\calH_0) \cap \calH_0$ is nontrivial. Indeed, each $\bbB^4_{r_j}$ meets some $\calD_k$ nontrivially, hence $R_j$ acts trivially on some totally geodesic subspace of $\calH_0$.

Since the $R_j$ generate $\PGam_4$ and $\PGam_4$ acts transitively on the set of short roots, it follows by induction that $\calH$ has connected support. Indeed, if
	\[
	\calH_d \coloneqq \bigcup_{|\gam| \le d}\gam(\calH_0)
	\]
for $|\gam|$ the word length in the generators $R_j$, then $\calH_0$ is connected for the base case. If $\calH_d$ is connected and $R_{j_1} \cdots R_{j_{d+1}}$ is a word of length $d+1$ in the generators $R_j$, then $R_{j_{d+1}}(\calH_0) \cap \calH_0 \neq \emptyset$, so
	\[
	(R_{j_1} \cdots R_{j_{d+1}})(\calH_0) \supseteq (R_{j_1} \cdots R_{j_d})\!\left(R_{j_{d+1}}(\calH_0) \cap \calH_0\right)\! \neq \emptyset
	\]
has nontrivial intersection with $(R_{j_1} \cdots R_{j_d})(\calH_0) \subseteq \calH_d$, which completes the induction.
\end{pf}

\begin{thebibliography}{99}
\bibitem[ACT]{ACT} D.~Allcock, J.~A. Carlson, and D.~Toledo. \newblock The complex hyperbolic geometry of the moduli space of cubic surfaces. \newblock {\em J. Algebr. Geom.}, 11(4):659--724, 2002.
\end{thebibliography}
\end{document}
